\documentclass[12pt,a4paper]{article}
\usepackage[margin=1in]{geometry}
\usepackage{setspace}
\usepackage{graphicx}
\usepackage{array}
\usepackage{longtable}
\usepackage{booktabs}
\usepackage{float}
\usepackage[hidelinks]{hyperref}

\setstretch{1.2}

\title{Assignment 1: Creating an AI-Powered Intrusion Detection Solution with Real-Time Threat Detection}
\author{Student Name: \underline{\hspace{6cm}}\\
Student ID: \underline{\hspace{6cm}}}
\date{Course: Information Security\\
Aligned CLO: 4 - Create solutions to real life scenarios using different security related tools}

\begin{document}
\maketitle
\thispagestyle{empty}
\newpage

\section*{Executive Summary}
This project presents a practical intrusion detection solution built for a SecureNet Corp. style scenario. I used the CICIDS2017 dataset to train and evaluate machine learning models for classifying network traffic as benign or malicious. The final runtime solution is a full stack system with a Django REST backend, an external API that streams sample traffic, and a Flutter dashboard that updates every 2 seconds. During local testing, both servers were healthy, model loading was successful, and a live test batch of 5 flows produced 2 threats and 3 benign records, with threat classes such as DoS Hulk and PortScan detected at high confidence. The dashboard also reports CIA+Au security impact for each threat type. The solution meets the assignment objective by showing a complete workflow from data preparation and model development to real-time deployment and validation.

\section{Introduction and Real-World Scenario}
Organizations receive a large volume of network traffic every day, and manual inspection cannot keep up with modern attack speed. Traditional controls such as firewalls are necessary, but they are not enough to classify advanced traffic patterns. In this project, I treated the task as a real request from a CISO: build a proof-of-concept tool that can classify network traffic and support security monitoring decisions.

The main goal was not only to train a model but also to operationalize it in a simple monitoring system. For that reason, I implemented a backend inference API and a live frontend dashboard. This makes the result closer to a real security operations workflow instead of only a notebook experiment.

\section{Methodology}
\subsection{Dataset and Scenario Justification}
I selected CICIDS2017 because it includes realistic and diverse attack behavior, including DDoS, DoS variants, brute-force attacks, port scanning, infiltration, and web attacks. These threats are relevant to enterprise web and network environments.

From the project pipeline and codebase:
\begin{itemize}
\item Dataset used for training: CICIDS2017 (full dataset loaded in the notebook pipeline).
\item Runtime simulation dataset: stratified sample of 500 rows for fast live testing.
\item Feature space: 78 flow-based features.
\item Label space: 15 classes (BENIGN + 14 attack classes).
\end{itemize}

\subsection{Data Preprocessing Steps}
The notebook workflow includes the following preprocessing sequence:
\begin{itemize}
\item CSV loading and merging with memory-conscious processing.
\item Label column detection and class distribution inspection.
\item Missing value and infinite value handling.
\item Label encoding with LabelEncoder.
\item Train/test split with stratification.
\item Standardization with StandardScaler.
\item Progressive class balancing (undersampling BENIGN and oversampling minority classes using SMOTE logic).
\end{itemize}

I also recorded an implementation detail that affected runtime stability: feature names in the scaler artifact required exact spacing from training time. The backend now uses an exact 78-column feature list and creates a pandas DataFrame before scaling to avoid inference mismatch errors.

\subsection{Machine Learning Tool Design}
The training notebook compares five model families:
\begin{itemize}
\item Logistic Regression
\item Random Forest
\item XGBoost
\item LightGBM
\item Neural Network
\end{itemize}

For deployment, XGBoost was chosen as the inference model in the backend service because it provided strong classification behavior with manageable inference cost for real-time polling.

The backend follows a singleton design for model loading, so model files are loaded once at startup and reused for each prediction call. This reduces repeated I/O and keeps latency low.

\section{Implemented Solution Architecture}
The final solution has three integrated parts:
\begin{enumerate}
\item Main backend (Django REST API): receives requests, classifies traffic, and returns statistics.
\item External API server (Django): provides traffic samples with API key validation.
\item Frontend dashboard (Flutter): polls every 2 seconds and visualizes statistics and threats.
\end{enumerate}

\subsection*{Main Backend Endpoints}
The backend exposes function-based REST endpoints:
\begin{itemize}
\item GET /api/health/
\item GET /api/statistics/
\item POST /api/predict/
\item GET /api/fetch-and-classify/
\item GET /api/attack-classes/
\item GET /api/recent-threats/
\item POST /api/reset-statistics/
\item GET /api/feature-names/
\end{itemize}

\subsection*{Real-Time Logic}
\begin{itemize}
\item Polling interval: 2 seconds (frontend provider timer).
\item Batch size per poll: 5 records (configurable through API query parameter).
\item Infinite traffic generation: external API samples with replacement from the 500-row runtime dataset.
\end{itemize}

\section{Results and Security Analysis}
\subsection{Current Local Runtime Test Evidence (19-Apr-2026)}
I performed direct API checks on localhost and recorded the following:
\begin{itemize}
\item External API health: healthy, data loaded, 500 records, 78 features.
\item Main backend health: healthy, model loaded successfully.
\item One fetch-and-classify run (batch size 5): 2 threats, 3 benign.
\item Updated statistics after run: total analyzed = 5, threat count = 2, benign count = 3, threat rate = 40.0\%.
\item Recent threat feed sample: PortScan and DoS Hulk with confidence around 0.999.
\end{itemize}

\subsection{CIA+Au Security Impact Mapping}
To improve security interpretation, each predicted class is mapped to:
\begin{itemize}
\item Confidentiality
\item Integrity
\item Availability
\item Authenticity
\end{itemize}

Example from runtime output:
\begin{itemize}
\item DoS Hulk impacts availability.
\item PortScan impacts confidentiality (reconnaissance exposure).
\end{itemize}

This mapping helps analysts move from raw class names to practical impact understanding.

\subsection{Critical Analysis}
The proof-of-concept meets the core assignment requirement and demonstrates practical detection in live polling mode. However, there are important limitations:
\begin{itemize}
\item Model artifact compatibility warnings appeared due to scikit-learn version mismatch between training and runtime.
\item Runtime currently uses a 500-row sampled source for fast testing, not full production-scale traffic.
\item API key is hardcoded for development convenience and should be moved to environment variables in production.
\end{itemize}

Security implication note: if recall is low for any attack class (especially stealth classes), the system may miss true attacks. This can create a false sense of safety. Therefore, confusion matrix review and class-wise recall tracking are mandatory before real deployment.

\section{Conclusion and Future Enhancements}
This project delivers a complete machine learning intrusion detection workflow aligned with CLO 4. It starts from realistic security data, applies preprocessing and model training, and ends with a working real-time dashboard.

In short, the solution is useful as a practical prototype and suitable for viva demonstration because all core components are integrated and testable.

Recommended next improvements:
\begin{itemize}
\item Retrain and export artifacts under pinned package versions to remove version mismatch risk.
\item Add class-wise threshold tuning for better recall on minority attack classes.
\item Replace hardcoded API key with environment-based secrets.
\item Add persistent storage for historical analytics and incident review.
\item Add automated test cases for endpoint correctness and model schema validation.
\end{itemize}

\section{Screenshot Placement Plan}
Paste screenshots in the report using the mapping below.

\begin{longtable}{p{0.12\linewidth} p{0.35\linewidth} p{0.45\linewidth}}
\toprule
ID & Paste In Section & What Screenshot Should Show \\
\midrule
S1 & Introduction / Methodology (Dataset) & Notebook cell showing CICIDS2017 files loaded and dataset shape summary. \\
S2 & Methodology (Preprocessing) & Notebook output showing cleaning, encoding, scaling, and class balancing steps. \\
S3 & Methodology (Model Design) & Notebook section listing models (LR, RF, XGBoost, LightGBM, NN). \\
S4 & Results and Security Analysis & Model comparison output (accuracy, precision, recall, F1, training time). \\
S5 & Results and Security Analysis & Confusion matrix figure for best model (or combined confusion matrix figure). \\
S6 & Implemented Architecture & Terminal screenshot: external API server running on 8001 with health success. \\
S7 & Implemented Architecture & Terminal screenshot: backend server running on 8080 with model loaded lines. \\
S8 & Results and Security Analysis & Terminal screenshot of /api/fetch-and-classify/ JSON response (threats + benign). \\
S9 & Results and Security Analysis & Terminal screenshot of /api/statistics/ response showing threat percentage. \\
S10 & Results and Security Analysis & Flutter dashboard main screen with cards and attack distribution chart. \\
S11 & Results and Security Analysis & Flutter recent threats panel showing CIA+Au badges/impact details. \\
S12 & Conclusion / Deployment Note & GitHub repository page after push (public repo with notebook and README visible). \\
\bottomrule
\end{longtable}

\subsection*{Suggested Caption Style}
Use short, direct captions such as:
\begin{itemize}
\item Figure S8: Live classification API response for batch size 5.
\item Figure S10: Real-time dashboard after monitoring started.
\item Figure S11: Recent threats with CIA+Au impact tags.
\end{itemize}

\section{References}
\begin{enumerate}
\item Canadian Institute for Cybersecurity. CICIDS2017 Dataset.
\item Scikit-learn documentation: preprocessing and evaluation metrics.
\item XGBoost documentation.
\item TensorFlow/Keras documentation.
\item Django and Django REST Framework documentation.
\item Flutter documentation.
\end{enumerate}

\section*{GitHub Profile and Repository}
GitHub Profile: \url{https://github.com/FahadIzMe}

Target Repository (public): \url{git@github.com:FahadIzMe/IS_1.git}

\end{document}
